\documentclass[11pt,a4paper]{article}
\usepackage[margin=1.8cm]{geometry}
\usepackage[T1]{fontenc}
\usepackage{lmodern}
\usepackage{booktabs}
\usepackage{tabularx}
\usepackage{array}
\usepackage[hidelinks]{hyperref}
\usepackage{xcolor}
\usepackage{enumitem}
\hypersetup{colorlinks=true, urlcolor=blue!60!black, linkcolor=black}

\newcolumntype{L}[1]{>{\raggedright\arraybackslash}p{#1}}
\newcolumntype{Y}{>{\raggedright\arraybackslash}X}
\renewcommand{\arraystretch}{1.3}
\setlength{\parindent}{0pt}
\setlength{\parskip}{0.5em}

% bullet list inside a table cell
\newenvironment{opts}
  {\begin{minipage}[t]{\linewidth}\raggedright\begin{itemize}[nosep,leftmargin=0.9em,label=\textbullet]}
  {\end{itemize}\vspace{3pt}\end{minipage}}

% literature links
\newcommand{\melcarne}{\href{https://www.eurecom.eu/en/publication/8539}{Melcarne}}
\newcommand{\sltnet}{\href{https://arxiv.org/abs/2412.12843v3}{SLTNet}}
\newcommand{\tevio}{\href{https://doi.org/10.1109/TIM.2025.3552392}{TEVIO}}
\newcommand{\openglass}{\href{https://arxiv.org/abs/2606.07431v1}{OpenGlass}}
\newcommand{\mcfnet}{\href{https://arxiv.org/abs/2508.10704}{MCFNet}}
\newcommand{\msad}{\href{https://openaccess.thecvf.com/content/CVPR2021/html/Qi_Multi-Scale_Aligned_Distillation_for_Low-Resolution_Detection_CVPR_2021_paper.html}{MSAD}}
\newcommand{\efinet}{\href{https://doi.org/10.1109/CVPRW53098.2021.00142}{EFI-Net}}
\newcommand{\li}{\href{https://doi.org/10.1609/aaai.v39i17.34013}{Li}}
\newcommand{\cafr}{\href{https://arxiv.org/abs/2407.12582}{CAFR}}
\newcommand{\thermalgyro}{\href{https://arxiv.org/abs/2506.12536}{Thermal-Gyro}}
\newcommand{\dagr}{\href{https://doi.org/10.1038/s41586-024-07409-w}{DAGR}}
\newcommand{\ramnet}{\href{https://doi.org/10.1109/LRA.2021.3060707}{RAMNet}}
\newcommand{\zhao}{\href{https://doi.org/10.1109/ICRA48891.2023.10160472}{Zhao}}
\newcommand{\mitrokhin}{\href{https://arxiv.org/abs/1803.04523}{Mitrokhin}}
\newcommand{\pedro}{\href{https://doi.org/10.1109/CVPRW59228.2023.00426}{PEDRo}}

\title{DalilSENSE Experimental Plan}
\begin{document}
\maketitle

\section*{Goal}
Show whether fusing $32\times24$ thermal with events, inside the network, labels each zone
(left / center / right) as clear, person, furniture or other better than the best single sensor,
at wearable-level compute.

\section*{Experiments}
{\footnotesize
\begin{tabularx}{\textwidth}{@{}L{0.6cm} L{3.1cm} Y L{2.2cm} L{2.4cm} L{1.0cm}@{}}
\toprule
\textbf{\#} & \textbf{Question} & \textbf{Compare} & \textbf{Keep fixed} & \textbf{Extra metrics} & \textbf{Owner} \\
\midrule
E1 & \textbf{Event representation:}\newline How to turn events into input: window length and form? &
  \begin{opts}
    \item Count frame over 40, 80, 120 or 250\,ms
    \item Fixed number of events
    \item Voxel grid (\melcarne, \sltnet)
    \item Time surface (\tevio)
    \item Polarity histogram (\openglass)
  \end{opts} &
  Events-only network & & \\
\addlinespace
E2 & \textbf{The resolution problem:}\newline How to join $32\times24$ thermal with $320\times320$ events? &
  \begin{opts}
    \item Shrink events at the input
    \item Learned downsampling (\mcfnet)
    \item Fuse where sizes match (\msad)
    \item Common low-res grid (\efinet)
  \end{opts} &
  \begin{opts}
    \item E1 input
    \item Concat fusion
  \end{opts} & MFLOPs & \\
\addlinespace
E3 & \textbf{Fusion method:}\newline Which fusion method (early, mid, late) gives the most gain for its compute? &
  \begin{opts}
    \item \textbf{Early:} stack inputs
    \item \textbf{Mid:} concat, add, scalar gate (\melcarne), channel gate (\li), attention (\cafr)
    \item \textbf{Late:} average, learned blend (\thermalgyro)
  \end{opts} &
  \begin{opts}
    \item E1 input
    \item E2 approach
  \end{opts} &
  F1 vs MFLOPs & \\
\addlinespace
E4 & \textbf{The frame-rate problem:}\newline One decision per thermal frame, or events updating between frames? &
  \begin{opts}
    \item Sync at 4\,Hz (now)
    \item Thermal anchor, events fill gaps (\dagr)
    \item Shared memory (\ramnet)
    \item Event clock (\tevio)
  \end{opts} &
  Best E3 model &
  \begin{opts}
    \item Decisions per second
    \item Latency
  \end{opts} & \\
\addlinespace
E5 & \textbf{The ego-motion problem:}\newline Does removing the wearer's own motion help? &
  \begin{opts}
    \item No compensation
    \item IMU de-rotation (\zhao)
    \item Compensation from events (\mitrokhin)
  \end{opts} &
  Best E3 model &
  \begin{opts}
    \item Walking straight vs turning
    \item Extra ms
  \end{opts} & \\
\addlinespace
E6 & \textbf{Classical baseline:}\newline Does a classical fusion work as well? &
  \begin{opts}
    \item Corridor event density (our test on \pedro)
    \item Hand-made features + logistic regression
  \end{opts} &
  E1 input & Cost & \\
\addlinespace
E7 & \textbf{Per-class results:}\newline Which classes gain: people, furniture or other objects? &
  Best single-sensor vs best fused (no retraining) &
  Trained models &
  \begin{opts}
    \item Per-class F1
    \item Confusion matrix
  \end{opts} & \\
\addlinespace
E8 & \textbf{Final comparison:}\newline \emph{Main:} does fusion beat the best single sensor at a reasonable cost? &
  Best of E1--E6 vs all baselines &
  5 seeds & All & \\
\bottomrule
\end{tabularx}
}

\textbf{Order:} E1 and E2 fix the inputs, E3 picks the fusion, E4 to E6 build on the best model, then E7 and E8 report.

\section*{Baselines (every experiment)}
{\small
\begin{tabularx}{\textwidth}{@{}L{6cm} Y@{}}
\toprule
\textbf{Baseline} & \textbf{Why} \\
\midrule

Thermal only, events only (one small CNN per sensor) & Single-sensor models to beat (about 0.1M parameters each) \\
Thermal only, double width & Rules out ``it's just a bigger model'' \\
Late average (each sensor answers alone, then the answers are averaged) & Simplest fusion \\
Oracle (trust whichever sensor is right) & The ceiling \\
\bottomrule
\end{tabularx}
}

\section*{Metrics (every experiment)}
{\small
\begin{tabularx}{\textwidth}{@{}L{6cm} Y@{}}
\toprule
\textbf{Metric} & \textbf{Why} \\
\midrule
\textbf{Per-zone macro F1 (main)} & Rare classes like furniture count equally \\
\% of zones correct & Easy to read \\
Per-class F1 & Which classes gain \\
Missed obstacles (zone called clear when it isn't) & False negatives are dangerous for navigation \\
Parameters, MFLOPs, CPU ms per update & Fits a wearable? \\
\bottomrule
\end{tabularx}
}

\section*{Setup (every experiment)}
{\small
\begin{tabularx}{\textwidth}{@{}L{3cm} Y@{}}
\toprule
\textbf{Item} & \textbf{Rule} \\
\midrule
Data & Our own wearable recordings: office, university, houses, supermarket; wearer always walking; RGB for labels only \\
Labels & YOLO on RGB, plus a hand-checked sample \\
Main split & Leave one place out; seen-place split secondary; no random frame splits \\
Seeds & 3 while exploring, 5 for the final run \\
Tuning & On a validation recording, never the test one \\
Fairness & Same backbone, recipe and folds for every model \\
``Better'' & Paired 95\% CI of the difference above zero \\
\bottomrule
\end{tabularx}
}

\section*{Open questions}
\begin{itemize}
  \item 
\end{itemize}

\end{document}
